\begin{table}[t]
\centering
\small
\setlength{\tabcolsep}{6pt}
\renewcommand{\arraystretch}{1.10}

\begin{tabular}{llccc}
\toprule
\textbf{Task} & \textbf{Model} & \textbf{Patient-level} & \textbf{Cell-level} & \textbf{Inflation} \\
\midrule

\multirow{2}{*}{Coarse lineage}
& XGBoost      & 0.322 & 0.372 & $+0.050$ \\
& \ourmethod{} & 0.333 & 0.432 & $+0.099$ \\
\midrule

\multirow{2}{*}{Fine subtype}
& XGBoost      & 0.143 & 0.212 & $+0.069$ \\
& \ourmethod{} & 0.156 & 0.261 & $+0.105$ \\

\bottomrule
\end{tabular}

\caption{
Macro-F1 under the primary patient-level 5-fold protocol versus the
cell-level protocol-sensitivity control (\S\ref{sec:data:protocols}), both
evaluated on the same models. Cell-level splitting, which allows cells from
the same patient to appear in both training and test sets, inflates
macro-F1 for both models, and inflates \ourmethod{}'s score by roughly
twice as much as XGBoost's. This indicates that \ourmethod{}'s
representation is more sensitive to patient-level leakage than a
tree-based baseline, underscoring why patient-independent evaluation is
necessary to avoid overstating hierarchical classification performance in
this small-cohort setting.
}
\label{tab:protocol_sensitivity}

\end{table}
